% Fig — how the two-variable heat PDE collapses to the one-variable closure
% ODE, in five moves. Mirrors app:pde2ode / sec:closure (equilibrium.py:
% _rectified_flux, _reconstruct_subskin). Body only; the book supplies
% \caption/\label.
\begin{tikzpicture}[
  font=\footnotesize,
  every node/.style={align=center},
  eqbox/.style={rectangle, rounded corners=2pt, draw=dim, line width=0.8pt,
                fill=white, inner sep=5pt, text width=46mm},
  keybox/.style={eqbox, draw=teal, fill=teal!7},
  res/.style={rectangle, rounded corners=3pt, draw=forest, fill=forest!12,
              line width=1pt, inner sep=5pt, text width=60mm},
  arr/.style={-{Stealth[length=2.2mm]}, line width=0.9pt, draw=dim,
              rounded corners=3pt, line cap=round},
  mv/.style={font=\scriptsize\bfseries\color{teal}}]

  \node[eqbox] (pde) at (0,0)
    {\textbf{the heat PDE} --- two variables $(z,t)$\\[2pt]
     $\rho c_p\,\partial_t T=\partial_z\!\left(K\,\partial_z T\right)$};

  \node[keybox] (avg) at (6.4,0)
    {\textbf{average over one lunation}\\[2pt]
     $\langle\,\cdot\,\rangle=\tfrac1P\!\int_0^P(\cdot)\,dt$ on a
     \emph{periodic} cycle\\[2pt]
     {\scriptsize\color{dim}$\langle\partial_t h\rangle
      =\tfrac1P[h(P)-h(0)]=0$: time drops out}};

  \node[eqbox] (const) at (12.8,0)
    {\textbf{flux constant in depth}\\[2pt]
     $\partial_z\langle K\,\partial_z T\rangle=0$
     $\Rightarrow$ $\langle K\,\partial_z T\rangle=Q_b$\\[2pt]
     {\scriptsize\color{dim}constant pinned by the basal BC}};

  \node[eqbox] (peel) at (12.8,-3.0)
    {\textbf{peel the average}\\[2pt]
     $u_{\rm rect}\equiv\langle K\partial_z T\rangle
      -K(\Tm)\,\partial_z\Tm$\\[2pt]
     {\scriptsize\color{dim}the oscillation--$K(T)$ correlation;\\
      $<\!1\%$ of $Q_b$ below the anchor}};

  \node[res] (ode) at (5.6,-3.0)
    {\textbf{the closure ODE} --- one variable $(z)$\\[2pt]
     $\dfrac{d\Tm}{dz}=\dfrac{Q_b-u_{\rm rect}}{K(\Tm,z)}$\\[2pt]
     {\scriptsize integrated downward from the anchor (Step B, RK2)}};

  \draw[arr] (pde) -- node[mv, above=1pt] {move 1--2} (avg);
  \draw[arr] (avg) -- node[mv, above=1pt] {move 3--4} (const);
  \draw[arr] (const) -- node[mv, right=2pt] {move 5} (peel);
  \draw[arr] (peel) -- (ode);
\end{tikzpicture}
